% q0030.tex — Leveraging the Situation
\question{
  id        = {0030},
  title     = {Leveraging the Situation},
  source    = {OBECON},
  year      = {2026},
  round     = {Seletiva IEO -- Phase B},
  language  = {English},
  topic     = {Finance},
  subtopic  = {Operating Leverage / Firm Theory},
  type      = {dissertative},
  statement = {%
    Consider two firms in different competitive markets, both facing the
    same price $P$. There is no entry or exit. Firm~A has fixed cost $FC_A$
    and marginal cost $MC_A(q) = 2 + q$. Firm~B has no fixed costs and the
    same marginal cost $MC_B(q) = 2 + q$.

    \begin{enumerate}[(a)]
      \item (20 points) For any given price $P$, determine the
            profit-maximising output for each firm.
            \textit{Hint: variable cost equals the area of the trapezoid
            under the marginal cost curve up to the optimal quantity.}

      \item (20 points) Suppose price doubles from $P = 6$ to $P' = 12$.
            Compute the absolute change in operating profit for Firm~A and
            Firm~B. Express answers as functions of $FC_A$ only.

      \item (15 points) Define the degree of operating leverage as
            \[
              \text{OL} = \frac{\%\,\Delta\,\text{operating profit}}
                               {\%\,\Delta\,\text{revenue}}.
            \]
            Calculate OL for both firms as price rises from $P=6$ to
            $P'=12$, and explain the implications.

      \item (25 points)
            \begin{enumerate}[(i)]
              \item Explain how operating leverage relates to variability of
                    a firm's profit.
              \item Name and briefly explain two other financial risks that
                    affect profit variability.
            \end{enumerate}

      \item (20 points) Provide two sectors with typically \textbf{high}
            operating leverage and two with \textbf{low} operating leverage,
            briefly justifying each.
    \end{enumerate}
  },
  choices   = {},
  answer    = {},
  solution  = {%
    \textbf{(a)} Variable cost: $VC(q) = 2q + q^2/2$. Shutdown condition:
    $P \ge AVC_{\min} = 2$. Profit-maximising output (if $P \ge 2$):
    $q^* = P - 2$.

    \textbf{(b)} General profit: $\Pi = (P-2)^2/2 - FC$.
    \begin{itemize}
      \item Firm A: $\Pi_A^t = 8 - FC_A$, $\Pi_A^{t+1} = 50 - FC_A$,
            $\Delta\Pi_A = 42$.
      \item Firm B: $\Pi_B^t = 8$, $\Pi_B^{t+1} = 50$,
            $\Delta\Pi_B = 42$.
    \end{itemize}
    Both have the same absolute change because $FC_A$ cancels.

    \textbf{(c)} Revenue at $P=6$: 24; at $P'=12$: 120.
    $\%\Delta R = 400\%$ for both.
    \[
      OL_A = \frac{42/(8-FC_A)}{4} = \frac{21}{16-2FC_A}, \qquad
      OL_B = \frac{42/8}{4} = \frac{21}{16} \approx 1.31.
    \]
    Since $FC_A > 0$ reduces the base profit $\Pi_A^t$, $OL_A > OL_B$:
    fixed costs amplify the sensitivity of profits to revenue changes.

    \textbf{(d.i)} Higher operating leverage amplifies both positive and
    negative revenue shocks into larger percentage profit swings, increasing
    earnings volatility.

    \textbf{(d.ii)} Financial leverage (fixed interest obligations amplify
    operating profit fluctuations into larger net-income swings) and
    foreign-exchange exposure (currency movements affect revenues or input
    costs independently of operations).

    \textbf{(e)} \textbf{High OL}: airlines (aircraft, infrastructure) and
    utilities (power plants, networks). \textbf{Low OL}: consulting (mainly
    variable labour) and staffing agencies (costs scale with placements).
  },
}
